% !TeX root = ../../Book1.tex
% 纲要标题：### 3.4 正多面体群
% 纲要位置：计划纲要.md，第 207 行。

\section{正多面体群}
\label{b1:sec:0304}

% 纲要内容（仅作写作计划，尚非教材正文）：
%
% * 正四面体、立方体与正八面体
% * 正二十面体与正十二面体的构造及对称
% * 顶点、边、面上的置换作用
% * 旋转群与含反射的全对称群的区别
%
% ---
%

几何图形的对称必须连同所在空间一起说明。正多边形在平面中的旋转与反射已经给出循环群和二面体群；进入三维以后，顶点、边、面和体对角线都可以作为被置换的对象。同一个空间旋转在这些集合上可能具有不同的置换符号。

\ProseParagraphBreak
以下把多面体中心放在原点，以保持图形的正交线性变换为对称。行列式为 \(1\) 的称为旋转，行列式为 \(-1\) 的反转空间定向。确定一个图形的对称群时，我们依次构造对称、证明没有其他对称，再通过具体同态识别群结构；本节的计数不预先借用下一章的群作用定理。

\begin{lemma}[三维旋转的固定轴]\label{b1:lem:rotation-axis}
设 \(T\) 是三阶实正交矩阵，且 \(\det T=1\)。若 \(T\ne I\)，则
\(\dim\ker(T-I)=1\)。这条固定直线的正交补在 \(T\) 下保持不变，\(T\) 在该平面上是一个非恒等的平面旋转。
\end{lemma}
\begin{proof}
由 \(T^{-1}=T^{\mathsf T}\) 及行列式的乘法性质，
\[
\begin{aligned}
\det(T-I)
&=\det T\det(I-T^{-1})
 =\det(I-T^{\mathsf T})\\
&=\det(I-T)=-\det(T-I).
\end{aligned}
\]
所以 \(\det(T-I)=0\)，存在非零固定向量 \(v\)。若 \(x\perp v\)，则
\(\langle Tx,v\rangle=\langle Tx,Tv\rangle=\langle x,v\rangle=0\)，故 \(v^\perp\) 保持不变。
在以 \(v/\lVert v\rVert\) 为首向量的正交基下，\(T\) 分成一个 \(1\) 与一个行列式为 \(1\) 的二阶正交块。
后一块的第一列可写成 \((\cos\theta,\sin\theta)^{\mathsf T}\)，第二列由正交性及定向唯一确定，故它是平面旋转矩阵。

若另有与 \(v\) 线性无关的固定向量，减去它在 \(v\) 上的投影，就得到 \(v^\perp\) 内的非零固定向量。
平面旋转固定一个非零向量便是恒等变换，于是整个 \(T=I\)，与假设矛盾。因此固定空间恰为一维，平面上的旋转也非恒等。
\end{proof}

为了从顶点坐标确定多面体，先说明本节所需的几个几何概念。有限点集 \(V=\{v_1,\ldots,v_m\}\) 的\term[tubao]{凸包}[convex hull]是所有凸组合
\[
\operatorname{conv}V
=\left\{\,\sum_{i=1}^m t_i v_i\ \middle|\ t_i\geq0,\ \sum_{i=1}^m t_i=1\,\right\}.
\]
若所有顶点满足 \(\langle a,v_i\rangle\leq b\)，其中 \(a\ne0\)，且至少一个顶点取等号，
则平面 \(\langle a,x\rangle=b\) 称为凸包的\term[zhichi-pingmian]{支持平面}[supporting plane]。
凸组合取等号时只能使用取等号的顶点，所以凸包与该平面的交恰是这些顶点的凸包。
交集若为二维多边形，便是多面体的一个面；例如恰有三个不共线顶点取等号时，交集是三角面。
指向 \(\langle a,x\rangle>b\) 一侧的法向称为该面的\term[wai-faxiang]{外法向}[outward normal]。

\ProseParagraphBreak
下面还会用截面判断面是否已经列全。用一个平面在某顶点附近截下一个小角，并让其余顶点都位于截面的另一侧，截面就是一个凸多边形。
经过原顶点的边给出截面的顶点，经过原顶点的面给出截面的边：在每个面内，截面切出连接相邻两条边的线段。
凸多边形等于各边所在直线的内侧闭半平面之交。因此，若几个已经找到的面在截面上围成一个多边形，截面既包含这个多边形，又被这些面的内侧半平面限制在其中，两者便相等。
原顶点与每条截面边唯一确定对应面的平面，故截面没有额外的边，就说明没有遗漏经过原顶点的面。

\subsection{正四面体、立方体与正八面体}
\label{b1:subsec:0304:01}

有限多面体的任意欧氏对称会置换顶点，所以保持所有顶点的平均值，也就是这里取作原点的中心。故把中心移到原点后，只需研究正交线性变换。它们的复合仍正交、逆仍正交，保持给定图形的那些变换形成其\term[jihe-duichengqun]{几何对称群}[symmetry group]，其中的旋转形成\term[xuanzhuanqun]{旋转群}[rotation group]。

\begin{proposition}[正四面体]\label{b1:prop:tetrahedral-group}
正四面体的全对称群在四个顶点上的置换给出与 \(S_4\) 的同构；旋转群对应 \(A_4\)。
\end{proposition}
\begin{proof}
取顶点
\begin{equation}\label{b1:eq:tetrahedron-vertices}
v_1=(1,1,1),\quad v_2=(1,-1,-1),\quad
v_3=(-1,1,-1),\quad v_4=(-1,-1,1).
\end{equation}
它们等长，任意两个不同顶点的内积为 \(-1\)，故构成正四面体。前三个向量线性无关，且 \(v_1+v_2+v_3+v_4=0\)。任意顶点排列都唯一延拓为线性变换：先规定前三个向量的像，和为零的关系自动规定第四个的像。由于基向量的两两内积不变，这个线性变换正交。因此每个顶点排列都实现为唯一几何对称。将对称送到其顶点排列的映射还保持复合，因为 \((TU)v_i=T(Uv_i)\)，故这个双射是群同构。

为判断定向，把四列 \((1,v_i)\) 排成一个四阶矩阵，其行列式非零。顶点排列 \(\sigma\) 使这个行列式乘以 \(\operatorname{sgn}(\sigma)\)；相应线性变换 \(T\) 则使它乘以 \(\det T\)。所以
\(\det T=\operatorname{sgn}(\sigma)\)，旋转恰对应偶置换。
\end{proof}

\begin{proposition}[立方体与正八面体]\label{b1:prop:cubic-rotation-group}
立方体与正八面体的旋转群同构于 \(S_4\)，各有 \(24\) 个元素；各自全对称群有 \(48\) 个元素。
\end{proposition}
\begin{proof}
取立方体顶点为 \((\pm1,\pm1,\pm1)\)。对称置换六个面的外法向
\(\pm e_1,\pm e_2,\pm e_3\)，所以矩阵的每列、每行恰有一个非零元，非零元为 \(\pm1\)。反过来，每个这样的带符号坐标置换都保持立方体。共有 \(3!\cdot2^3=48\) 个，其中行列式正负各半，故旋转共 \(24\) 个。同一组六个法向本身构成正八面体的顶点集，因此两个图形有相同的正交对称群。

立方体的四条体对角线可由\eqref{b1:eq:tetrahedron-vertices}中的 \(v_1,v_2,v_3,v_4\) 表示，每条线忽略向量的正负。旋转置换这四条线，给出到 \(S_4\) 的同态。若旋转 \(T\) 固定每条线，则
\(Tv_i=\varepsilon_i v_i\)，\(\varepsilon_i\in\{1,-1\}\)。这四向量的唯一线性关系由和为零生成，所以
\(\sum_i\varepsilon_i v_i=0\) 迫使四个 \(\varepsilon_i\) 相同。故
\(T=I\) 或 \(T=-I\)；后一矩阵在三维中行列式为 \(-1\)，不是旋转。核因而平凡，\cref{b1:prop:kernel-fibers}保证这个同态单射；两边都有 \(24\) 个元素，所以是同构。
\end{proof}

\subsection{正二十面体与正十二面体}
\label{b1:subsec:0304:icosahedron-dodecahedron}

正二十面体也可以从坐标构造。这里先确认顶点、边和面的完整结构，再计算对称群；正十二面体则由这些面的中心构造。令
\(\phi=(1+\sqrt5)/2\)，取十二个顶点
\[
V=\bigl\{(0,\pm1,\pm\phi),\ (\pm1,\pm\phi,0),\
(\pm\phi,0,\pm1)\bigr\}.
\]
记 \(P=\operatorname{conv}V\)。下面验证它是正二十面体。
各顶点的长度的平方为 \(\phi+2\)，最短的顶点间距离为 \(2\)。以
\(v=(0,1,\phi)\) 为例，距离为 \(2\) 的五个点依次记为 \(w_1,\ldots,w_5\)：
\[
(0,-1,\phi),\ (\phi,0,1),\ (1,\phi,0),\
(-1,\phi,0),\ (-\phi,0,1).
\]
它们与 \(v\) 的内积都为 \(\phi\)，在垂直于 \(v\) 的同一平面中构成正五边形：相邻点距离为 \(2\)，非相邻点距离为 \(2\phi\)。
以下下标按模 \(5\) 计算。五个三角形 \(\operatorname{conv}\{v,w_i,w_{i+1}\}\) 都是等边三角形。
对于其中任意一个，记其三个顶点为 \(u,v,w\)，则
\[
\langle u+v+w,u\rangle
=\langle u+v+w,v\rangle
=\langle u+v+w,w\rangle=3\phi+2.
\]
不同顶点间的内积至多为 \(\phi\)，故对其余顶点 \(z\)，
\(\langle u+v+w,z\rangle\leq3\phi\)。因此每个候选三角形确实是 \(P\) 的面。

\ProseParagraphBreak
为排除遗漏，在 \(v\) 附近取截面
\[
\langle v,x\rangle=\phi+2-\varepsilon,
\qquad 0<\varepsilon<2.
\]
除 \(v\) 外的所有顶点在此平面的另一侧。它与 \([v,w_i]\) 的交点为
\[
p_i=\left(1-\frac{\varepsilon}{2}\right)v
     +\frac{\varepsilon}{2}w_i.
\]
五个已证三角面与截面相交成五条线段 \([p_i,p_{i+1}]\)，它们围成一个正五边形。
截面包含这五个点的凸包，又位于这五条边的各自内侧，所以恰是这个正五边形。
经过 \(v\) 的边和面因而各有五个，正是已经列出的那些。

\ProseParagraphBreak
坐标循环置换和任意两个坐标同时变号保持 \(V\)，并能把 \(v\) 送到任意顶点，所以每个顶点都有同样的邻接结构。
每条边有两个端点，每个三角面有三个顶点，故
\[
E=\frac{12\cdot5}{2}=30,
\qquad F=\frac{12\cdot5}{3}=20.
\]
这些面全是全等的等边三角形，每个顶点恰有五个面相交，得到正二十面体。

\begin{proposition}[正二十面体与正十二面体]\label{b1:prop:icosahedral-order}
正二十面体与正十二面体各有 \(60\) 个旋转对称，\(120\) 个全对称。
\end{proposition}
\begin{proof}
坐标循环置换与任意两个坐标同时变号都是行列式为 \(1\) 的对称，它们能把上面的 \(v\) 送到十二个顶点中的任意一个。固定 \(v\) 的旋转必须置换其五个邻点，并在垂直于 \(v\) 的平面中保持定向，故只能是该正五边形的五个旋转之一。这五个旋转确实保持整个顶点集：它们固定 \(v,-v\)，置换五个邻点，也置换这五点的相反点，而这些点恰好构成全部十二个顶点。

对每个目标顶点 \(w\)，固定一个把 \(v\) 送到 \(w\) 的旋转 \(t_w\)。所有把 \(v\) 送到 \(w\) 的旋转恰为 \(t_w\) 与那五个固定 \(v\) 的旋转的复合：左乘 \(t_w^{-1}\) 即可验证这一一对应。所以总数为 \(12\cdot5=60\)。中心反演 \(-I\) 保持顶点集且反转定向，与它复合建立两类定向之间的双射，故全对称群有 \(120\) 个元素。

再构造正十二面体。对每个三角面 \(F\) 的三个顶点 \(u,w,z\)，令
\[
c_F=\frac{u+w+z}{3},\qquad
Q=\operatorname{conv}\{\,c_F\mid F\text{ 是 }P\text{ 的三角面}\,\}.
\]
这里 \(c_F\) 是等边三角面的中心。所有这些中心满足
\(\lVert c_F\rVert^2=\phi+\frac23\)。这些面中心组成的点集关于原点对称。
固定原顶点 \(v\)，若 \(v\in F\)，则
\[
\langle v,c_F\rangle
=\frac{(\phi+2)+\phi+\phi}{3}=\phi+\frac23;
\]
若 \(v\notin F\)，则 \(\langle v,c_F\rangle\leq\phi\)。特别地，若 \(F\ne F'\)，取 \(v\in F\setminus F'\)，这两种内积取值便说明 \(c_F\ne c_{F'}\)。
因此，经过 \(v\) 的五个三角面的中心恰在 \(Q\) 的支持平面
\(\langle v,x\rangle=\phi+\frac23\) 上。
绕 \(v\) 轴的五次旋转循环置换这五点，所以它们围成正五边形面，记为 \(Q_v\)。
坐标对称把各个 \(v\) 互相变换，故这十二个五边形全等。
一个五边形及其相反面不在同一平面内，故 \(Q\) 是三维凸多面体。

还须证明 \(Q\) 没有其他面。若原来的边 \([u,w]\) 连接两三角面 \(F,F'\)，则同时满足
\(\langle u,x\rangle=\langle w,x\rangle=\phi+\frac23\)
的面中心恰为 \(c_F,c_{F'}\)。因此
\(Q_u\cap Q_w=[c_F,c_{F'}]\)，这是两个五边形共有的边。
在一个固定的 \(c_F\) 处，原三角面 \(F\) 的三条边给出这样的三条边，它们两两同属 \(F\) 的一个顶点所对应的五边形。

由于各面中心等长且两两不同，线性函数 \(x\mapsto\langle c_F,x\rangle\) 在这些点中仅于 \(c_F\) 取最大值，故每个 \(c_F\) 确为 \(Q\) 的顶点。
取 \(\eta>0\) 小于所有正数
\(\lVert c_F\rVert^2-\langle c_F,c_{F'}\rangle\)，其中 \(F'\ne F\)，用
\(\langle c_F,x\rangle=\lVert c_F\rVert^2-\eta\) 作截面。
上述三条边给出三个截点。每一对截点分别位于同一个正五边形在 \(c_F\) 处的两条不同邻边上，故两两不同。
它们也不共线：否则这条直线同时位于两个五边形的平面内，只能是它们的公共边所在直线；但那条直线经过 \(c_F\)，与截面只有一个交点。
所以三个已证五边形在截面上给出一个三角形的三条边。
截面包含这个三角形，又被这三条边的内侧半平面限制在其中，故恰是三角形。
所以每个 \(c_F\) 处只有这三个面；任何面都含有顶点，十二个 \(Q_v\) 因而已经列全。
新多面体有 \(20\) 个顶点、\(20\cdot3/2=30\) 条边及 \(12\) 个全等正五边形面，即正十二面体。

保持 \(P\) 的正交变换置换其三角面，也就置换面中心并保持 \(Q\)。
反过来，保持 \(Q\) 的正交变换置换已列全的十二个面，其单位外法向恰为
\(v/\sqrt{\phi+2}\)，其中 \(v\in V\)。因此它也保持 \(V\)，从而保持 \(P\)。
两图形的正交对称群相同，旋转群也相同，所述群阶随之得到。
\end{proof}

这里直接求出了两种群的阶，并没有把阶相同当作同构证明。识别正二十面体旋转群与 \(A_5\) 的进一步关系，需要另选合适的五个几何对象或补充群结构论证，本节不使用这一识别作为前提。

\subsection{顶点、边与面上的置换作用}
\label{b1:subsec:0304:02}

每个几何对称都置换顶点、边和面，且复合对称所得的置换是原置换的复合。因此选择这些有限集合，就得到从几何对称群到某个对称群的同态。这是把连续空间中的几何运动化为有限代数计算的方法；下一章将把这种做法统一称为群作用。

\ProseParagraphBreak
对本节的正多面体，顶点置换能够区分全部对称：固定全部顶点的线性变换固定一组空间基，因而是恒等变换。边置换也能区分对称，因为一个顶点是所有与它相接的边的唯一公共点；若每条边作为集合均被固定，每个顶点也被固定。类似地，一个顶点由与它相接的面的公共交点确定，所以固定每个面的对称也是恒等对称。

\ProseParagraphBreak
同一对称在不同集合上可以有很不一样的轮换分解。例如立方体绕 \(z\) 轴逆时针转过 \(\pi/2\)，映射为
\(T(x,y,z)=(-y,x,z)\)。在八个顶点上，它分别循环置换上面与下面的四个顶点，轮换类型为 \(4+4\)，因此是偶置换。在六个面上，上下面固定，四个侧面循环，轮换类型为 \(4+1+1\)，因此是奇置换。在四条体对角线上，按\eqref{b1:eq:tetrahedron-vertices}中的 \(v_i\) 编号，它给出 \((1\,2\,4\,3)\)，也是奇置换。

\ProseParagraphBreak
这说明“空间旋转对应偶置换”并非普遍规则。对于四面体顶点它成立，是因为已经证明三维行列式等于该顶点排列的符号；换成另一套被置换对象后，必须重新判断。前面的立方体旋转群同构 \(S_4\)，也正是四条体对角线的置换所给出的，而非八个顶点的全部置换群。

\begin{proposition}\label{b1:prop:cube-lines-kernel}
立方体全对称群在四条体对角线上的置换同态的核为 \(\{I,-I\}\)，而在八个顶点上的置换同态的核为 \(\{I\}\)。
\end{proposition}
\begin{proof}
设 \(T\) 固定四条体对角线。沿用\eqref{b1:eq:tetrahedron-vertices}的四个等长向量，正交性给出 \(Tv_i=\varepsilon_i v_i\)，其中 \(\varepsilon_i\in\{1,-1\}\)。由于 \(v_1,v_2,v_3\) 线性无关且 \(v_4=-v_1-v_2-v_3\)，等式 \(\sum_i Tv_i=0\) 变为
\[
(\varepsilon_1-\varepsilon_4)v_1+
(\varepsilon_2-\varepsilon_4)v_2+
(\varepsilon_3-\varepsilon_4)v_3=0.
\]
所以四个符号相同，\(T=I\) 或 \(T=-I\)。这重用了\cref{b1:prop:cubic-rotation-group}证明中的线性关系，但此处允许反转定向，所以两者都保留；它们也确实固定每条体对角线。若顶点置换恒等，则 \(T\) 固定基向量 \(v_1,v_2,v_3\)，只能是 \(I\)。
\end{proof}

所以顶点、边、面或对角线的选择会改变同态的核。较小的置换集合通常使计算更简洁，但必须检查它是否仍能分辨全部对称，或者明确记录哪些对称被视为相同。

\subsection{旋转群与全对称群}
\label{b1:subsec:0304:03}

行列式是从全对称群到 \(\{1,-1\}\) 的同态，旋转群恰为其核。只要存在一个反转定向的对称，该行列式同态就满射到二元素群；由\cref{b1:thm:first-isomorphism}，全对称群商去旋转群同构于 \(\{1,-1\}\)，故旋转群的指数为 \(2\)。五种正多面体都有这样的对称：四面体由奇顶点排列实现，其余四种由中心反演实现。因此全对称群的元素数都是旋转群的两倍。\cref{b1:tab:regular-polyhedra}汇总了本节的计算。

\begin{table}[htbp]
\centering
\caption[正多面体的对称计数]{正多面体的对称计数。旋转与全对称均指三维空间中的等距对称。}
\label{b1:tab:regular-polyhedra}
\begin{tabular}{lrrrrr}
\toprule
多面体 & 顶点数 & 边数 & 面数 & 旋转群阶 & 全对称群阶 \\
\midrule
正四面体 & 4 & 6 & 4 & 12 & 24 \\
立方体 & 8 & 12 & 6 & 24 & 48 \\
正八面体 & 6 & 12 & 8 & 24 & 48 \\
正十二面体 & 20 & 30 & 12 & 60 & 120 \\
正二十面体 & 12 & 30 & 20 & 60 & 120 \\
\bottomrule
\end{tabular}
\end{table}

其中正四面体、立方体和正八面体的旋转群已分别识别为 \(A_4,S_4,S_4\)；正十二面体与正二十面体在本节完成了对称群阶数的计算及两者之间的对应，旋转群与 \(A_5\) 的同构尚未在这里证明。

\ProseParagraphBreak
中心对称的三维图形还有一个更具体的分解。若 \(R\) 是其旋转群，则任意全对称都唯一写成 \((-I)^\varepsilon r\)，其中 \(r\in R\)、\(\varepsilon\in\{0,1\}\)。存在性由行列式决定是否乘一次 \(-I\)，唯一性来自 \(-I\) 不是旋转；又因 \(-I\) 与所有线性变换交换，乘法为
\[
\bigl((-I)^\varepsilon r\bigr)\bigl((-I)^\delta s\bigr)
=(-I)^{\varepsilon+\delta}(rs).
\]
因此立方体、正八面体、正十二面体、正二十面体的全对称群都可看作一个旋转及一个二值定向标记组成的群，乘法逐项进行。后面讨论群的直积时将给这种结构统一命名。

\ProseParagraphBreak
四面体没有中心反演对称：\(-v_1\) 不是它的顶点。因此它的全对称群虽也比旋转群大一倍，却不能直接沿用上面的分解。实际上全对称群 \(S_4\) 的中心平凡：若某置换与每个换位交换，\eqref{b1:eq:cycle-conjugation}说明它固定每个二元素子集。对不同的 \(i,j,k\)，单点集 \(\{i\}=\{i,j\}\cap\{i,k\}\) 因而也被固定，故它固定每个点。若四面体全对称群也有那样独立且可交换的二值标记，就会出现一个非平凡中心元素，与此矛盾。相同的阶数比，并不意味着相同的扩张方式。
